\documentclass[10pt,a4paper]{amsart}
\usepackage{amsmath,amssymb,amsthm,graphicx}
% Retained minimal preamble for the Stacks Project chapter more-algebra.tex (More on Algebra).
% The chapter file pinned-original-source.tex carries only \input{preamble} on its first line, so no
% preamble text lives inside the chapter itself.  The Stacks Project uses one single shared preamble for
% every chapter; the snapshot stages it identically as sources/stacks-chow/preamble.tex and
% sources/stacks-etale-cohomology/preamble.tex.  The block below is copied mechanically, line for line,
% from sources/stacks-chow/preamble.tex lines 14-18 (the Xy-pic setup) and 154-217 (the native theorem
% declarations together with the mathematical macros).  Nothing was added, removed or redefined; the
% wrapper class and the amsmath/amssymb/amsthm/graphicx base packages are supplied by the assembler.
% Chapter: more-algebra.tex -- More on Algebra.
% For commutative diagrams we use Xy-pic
\usepackage[all]{xy}

% We use 2cell for 2-commutative diagrams.
\xyoption{2cell}

% Theorem environments.
%
\theoremstyle{plain}
\newtheorem{theorem}[subsection]{Theorem}
\newtheorem{proposition}[subsection]{Proposition}
\newtheorem{lemma}[subsection]{Lemma}

\theoremstyle{definition}
\newtheorem{definition}[subsection]{Definition}
\newtheorem{example}[subsection]{Example}
\newtheorem{exercise}[subsection]{Exercise}
\newtheorem{situation}[subsection]{Situation}

\theoremstyle{remark}
\newtheorem{remark}[subsection]{Remark}
\newtheorem{remarks}[subsection]{Remarks}

\numberwithin{equation}{subsection}

% Macros
%
\def\lim{\mathop{\mathrm{lim}}\nolimits}
\def\colim{\mathop{\mathrm{colim}}\nolimits}
\def\Spec{\mathop{\mathrm{Spec}}}
\def\Hom{\mathop{\mathrm{Hom}}\nolimits}
\def\Ext{\mathop{\mathrm{Ext}}\nolimits}
\def\SheafHom{\mathop{\mathcal{H}\!\mathit{om}}\nolimits}
\def\SheafExt{\mathop{\mathcal{E}\!\mathit{xt}}\nolimits}
\def\Sch{\mathit{Sch}}
\def\Mor{\mathop{\mathrm{Mor}}\nolimits}
\def\Ob{\mathop{\mathrm{Ob}}\nolimits}
\def\Sh{\mathop{\mathit{Sh}}\nolimits}
\def\NL{\mathop{N\!L}\nolimits}
\def\CH{\mathop{\mathrm{CH}}\nolimits}
\def\proetale{{pro\text{-}\acute{e}tale}}
\def\etale{{\acute{e}tale}}
\def\QCoh{\mathit{QCoh}}
\def\Ker{\mathop{\mathrm{Ker}}}
\def\Im{\mathop{\mathrm{Im}}}
\def\Coker{\mathop{\mathrm{Coker}}}
\def\Coim{\mathop{\mathrm{Coim}}}

% Boxtimes
%
\DeclareMathSymbol{\boxtimes}{\mathbin}{AMSa}{"02}

%
% Macros for moduli stacks/spaces
%
\def\QCohstack{\mathcal{QC}\!\mathit{oh}}
\def\Cohstack{\mathcal{C}\!\mathit{oh}}
\def\Spacesstack{\mathcal{S}\!\mathit{paces}}
\def\Quotfunctor{\mathrm{Quot}}
\def\Hilbfunctor{\mathrm{Hilb}}
\def\Curvesstack{\mathcal{C}\!\mathit{urves}}
\def\Polarizedstack{\mathcal{P}\!\mathit{olarized}}
\def\Complexesstack{\mathcal{C}\!\mathit{omplexes}}
% \Pic is the operator that assigns to X its picard group, usage \Pic(X)
% \Picardstack_{X/B} denotes the Picard stack of X over B
% \Picardfunctor_{X/B} denotes the Picard functor of X over B
\def\Pic{\mathop{\mathrm{Pic}}\nolimits}
\def\Picardstack{\mathcal{P}\!\mathit{ic}}
\def\Picardfunctor{\mathrm{Pic}}
\def\Deformationcategory{\mathcal{D}\!\mathit{ef}}

% --- appended 2026-10-12 (Chief): the project own definitions of the reference/slogan/history
% display environments, copied mechanically from the shared Stacks preamble lines 8-12
% (\usepackage{verbatim} plus the three \newenvironment lines, which the chapter preamble keeps
% verbatim from the project preamble). Nothing added, removed or redefined: these environments
% carry no mathematical content and the project itself typesets them as comments. Without these
% three lines a counted statement whose own text uses them fails with Environment slogan undefined.
\usepackage{verbatim}
\newenvironment{reference}{\comment}{\endcomment}
\newenvironment{slogan}{\comment}{\endcomment}
\newenvironment{history}{\comment}{\endcomment}
\usepackage{cleveref}
\crefname{originallemma}{original lemma}{original lemmas}
\Crefname{originallemma}{Original lemma}{Original lemmas}
\makeatletter
\newlabel{more-algebra-lemma-fibres-henselization}{{07QQ}{source}{}{sourcepointer-1}{}}
\newlabel{more-algebra-lemma-fibres-henselization@cref}{{[originallemma][0][]07QQ}{[1][0][]source}{}{}{}}
\makeatother
\crefname{originallemma}{original lemma}{original lemmas}
\Crefname{originallemma}{Original lemma}{Original lemmas}
\makeatletter
\newlabel{lemma-filtered-colimit-etale-noetherian-fibres}{{0AH1}{source}{}{sourcepointer-2}{}}
\newlabel{lemma-filtered-colimit-etale-noetherian-fibres@cref}{{[originallemma][0][]0AH1}{[1][0][]source}{}{}{}}
\makeatother
\crefname{originallemma}{original lemma}{original lemmas}
\Crefname{originallemma}{Original lemma}{Original lemmas}
\makeatletter
\newlabel{lemma-dumb-properties-henselization}{{07QM}{source}{}{sourcepointer-3}{}}
\newlabel{lemma-dumb-properties-henselization@cref}{{[originallemma][0][]07QM}{[1][0][]source}{}{}{}}
\makeatother
\crefname{originallemma}{original lemma}{original lemmas}
\Crefname{originallemma}{Original lemma}{Original lemmas}
\makeatletter
\newlabel{lemma-henselization-noetherian}{{06LJ}{source}{}{sourcepointer-4}{}}
\newlabel{lemma-henselization-noetherian@cref}{{[originallemma][0][]06LJ}{[1][0][]source}{}{}{}}
\makeatother
\crefname{originallemma}{original lemma}{original lemmas}
\Crefname{originallemma}{Original lemma}{Original lemmas}
\makeatletter
\newlabel{algebra-lemma-quotient-henselization}{{05WQ}{source}{}{sourcepointer-5}{}}
\newlabel{algebra-lemma-quotient-henselization@cref}{{[originallemma][0][]05WQ}{[1][0][]source}{}{}{}}
\makeatother
\crefname{originallemma}{original lemma}{original lemmas}
\Crefname{originallemma}{Original lemma}{Original lemmas}
\makeatletter
\newlabel{algebra-lemma-quotient-strict-henselization}{{05WS}{source}{}{sourcepointer-6}{}}
\newlabel{algebra-lemma-quotient-strict-henselization@cref}{{[originallemma][0][]05WS}{[1][0][]source}{}{}{}}
\makeatother
\crefname{originallemma}{original lemma}{original lemmas}
\Crefname{originallemma}{Original lemma}{Original lemmas}
\makeatletter
\newlabel{lemma-henselization-reduced}{{06DH}{source}{}{sourcepointer-7}{}}
\newlabel{lemma-henselization-reduced@cref}{{[originallemma][0][]06DH}{[1][0][]source}{}{}{}}
\makeatother
\crefname{originallemma}{original lemma}{original lemmas}
\Crefname{originallemma}{Original lemma}{Original lemmas}
\makeatletter
\newlabel{algebra-proposition-dimension-zero-ring}{{00KJ}{source}{}{sourcepointer-8}{}}
\newlabel{algebra-proposition-dimension-zero-ring@cref}{{[originallemma][0][]00KJ}{[1][0][]source}{}{}{}}
\makeatother
\crefname{originallemma}{original lemma}{original lemmas}
\Crefname{originallemma}{Original lemma}{Original lemmas}
\makeatletter
\newlabel{algebra-lemma-minimal-prime-reduced-ring}{{00EU}{source}{}{sourcepointer-9}{}}
\newlabel{algebra-lemma-minimal-prime-reduced-ring@cref}{{[originallemma][0][]00EU}{[1][0][]source}{}{}{}}
\makeatother
\crefname{originallemma}{original lemma}{original lemmas}
\Crefname{originallemma}{Original lemma}{Original lemmas}
\makeatletter
\newlabel{algebra-lemma-etale-at-prime}{{00U4}{source}{}{sourcepointer-10}{}}
\newlabel{algebra-lemma-etale-at-prime@cref}{{[originallemma][0][]00U4}{[1][0][]source}{}{}{}}
\makeatother
\title{The fibres of the henselization and strict henselization of a local ring}
\author{The Stacks Project Authors}
\date{stacks-project:more-algebra.tex:local-snapshot-20261012; GNU FDL 1.2 or later}
\begin{document}
\maketitle
\setcounter{subsection}{1}
\noindent\textit{Editorial note.} All delivered mathematical text is the author's own, copied from the pinned Stacks Project chapter file more-algebra.tex; no statement or proof step was rewritten, reordered or removed.\par\medskip
\noindent\textit{Editorial note.} The chapter file pinned-original-source.tex carries only \textbackslash{}input\{preamble\} on its first line, so no preamble text lives inside the chapter. The retained file sources/stacks-more-algebra/retained-preamble.tex is a mechanical, line-for-line copy of the project's single shared preamble as staged in the snapshot (Xy-pic setup plus theorem declarations and mathematical macros). Nothing was added, removed or redefined, and the excerpt serves only as preamble, never as a span coordinate.\par\medskip
\noindent\textit{Editorial note.} Both counted spans start at the \textbackslash{}begin line of their environment, so the statement and the author's proof are delivered whole; surrounding prose of the section is deliberately outside the counted spans and is carried as standard prerequisites instead.\par\medskip
\noindent\textit{Editorial note.} The citation fields (source\_id, source\_version, authors, source\_url, license.name, source\_title) are identical to those of every other card this lane has produced from this chapter, so the batch selector treats them as one source.\par\medskip
\noindent\textit{Allowed standard prerequisites (editorial).}\begin{itemize}
\item Basic properties of henselization, and Noetherianity of the henselization and strict henselization. Lemmas \textbackslash{}ref\{lemma-dumb-properties-henselization\} and \textbackslash{}ref\{lemma-henselization-noetherian\}, cited at proof lines 11295-11296. Supply the properties used without mention throughout, in particular finiteness of the number of minimal primes.
\item The henselization of a quotient is the quotient of the henselization. Algebra, Lemma \textbackslash{}ref\{algebra-lemma-quotient-henselization\} and Algebra, Lemma \textbackslash{}ref\{algebra-lemma-quotient-strict-henselization\}, cited at proof lines 11298-11300. Allow the reduction to the case where the ring is a domain and the prime is zero.
\item A domain has reduced henselization. Lemma \textbackslash{}ref\{lemma-henselization-reduced\}, cited at proof line 11307. Makes the base change reduced, so that its primes are minimal and maximal at once.
\item A reduced Noetherian ring of dimension zero is a product of fields. Algebra, Proposition \textbackslash{}ref\{algebra-proposition-dimension-zero-ring\} and Algebra, Lemma \textbackslash{}ref\{algebra-lemma-minimal-prime-reduced-ring\}, cited at proof lines 11314-11317. Identify the base change with the product of the residue fields of the primes over the prime.
\item Residue field extensions of etale ring maps are finite separable. Algebra, Lemma \textbackslash{}ref\{algebra-lemma-etale-at-prime\}, cited at proof lines 11321-11324. Shows that the residue field extensions obtained from the colimit presentation are separable algebraic.
\end{itemize}\medskip
\ifdefined\hypertarget\hypertarget{sourcepointer-1}{}\fi
\noindent\textit{Source reference.} Original lemma 07QQ: focused-seeds/AGP/sources/stacks-snapshot-20261011/tags.txt, source lines 9840--9840. The lemma's own canonical tag, quoted so that readers can resolve the label inside the delivered text; it points at the counted statement itself and carries no proof content.\par\medskip
\ifdefined\hypertarget\hypertarget{sourcepointer-2}{}\fi
\noindent\textit{Source reference.} Original lemma 0AH1: focused-seeds/AGP/sources/stacks-snapshot-20261011/tags.txt, source lines 13214--13214. Navigation pointer to a statement the counted proof refers to but does not contain; nothing load-bearing is delegated to it.\par\medskip
\ifdefined\hypertarget\hypertarget{sourcepointer-3}{}\fi
\noindent\textit{Source reference.} Original lemma 07QM: focused-seeds/AGP/sources/stacks-snapshot-20261011/tags.txt, source lines 9837--9837. Navigation pointer to a statement the counted proof refers to but does not contain; nothing load-bearing is delegated to it.\par\medskip
\ifdefined\hypertarget\hypertarget{sourcepointer-4}{}\fi
\noindent\textit{Source reference.} Original lemma 06LJ: focused-seeds/AGP/sources/stacks-snapshot-20261011/tags.txt, source lines 8468--8468. Navigation pointer to a statement the counted proof refers to but does not contain; nothing load-bearing is delegated to it.\par\medskip
\ifdefined\hypertarget\hypertarget{sourcepointer-5}{}\fi
\noindent\textit{Source reference.} Original lemma 05WQ: focused-seeds/AGP/sources/stacks-snapshot-20261011/tags.txt, source lines 7594--7594. Navigation pointer to a statement the counted proof refers to but does not contain; nothing load-bearing is delegated to it.\par\medskip
\ifdefined\hypertarget\hypertarget{sourcepointer-6}{}\fi
\noindent\textit{Source reference.} Original lemma 05WS: focused-seeds/AGP/sources/stacks-snapshot-20261011/tags.txt, source lines 7596--7596. Navigation pointer to a statement the counted proof refers to but does not contain; nothing load-bearing is delegated to it.\par\medskip
\ifdefined\hypertarget\hypertarget{sourcepointer-7}{}\fi
\noindent\textit{Source reference.} Original lemma 06DH: focused-seeds/AGP/sources/stacks-snapshot-20261011/tags.txt, source lines 8181--8181. Navigation pointer to a statement the counted proof refers to but does not contain; nothing load-bearing is delegated to it.\par\medskip
\ifdefined\hypertarget\hypertarget{sourcepointer-8}{}\fi
\noindent\textit{Source reference.} Original lemma 00KJ: focused-seeds/AGP/sources/stacks-snapshot-20261011/tags.txt, source lines 753--753. Navigation pointer to a statement the counted proof refers to but does not contain; nothing load-bearing is delegated to it.\par\medskip
\ifdefined\hypertarget\hypertarget{sourcepointer-9}{}\fi
\noindent\textit{Source reference.} Original lemma 00EU: focused-seeds/AGP/sources/stacks-snapshot-20261011/tags.txt, source lines 548--548. Navigation pointer to a statement the counted proof refers to but does not contain; nothing load-bearing is delegated to it.\par\medskip
\ifdefined\hypertarget\hypertarget{sourcepointer-10}{}\fi
\noindent\textit{Source reference.} Original lemma 00U4: focused-seeds/AGP/sources/stacks-snapshot-20261011/tags.txt, source lines 1098--1098. Navigation pointer to a statement the counted proof refers to but does not contain; nothing load-bearing is delegated to it.\par\medskip
\begin{lemma}
\label{lemma-fibres-henselization}
Let $R$ be a Noetherian local ring. Let $\mathfrak p \subset R$ be a prime.
Then
$$
R^h \otimes_R \kappa(\mathfrak p) =
\prod\nolimits_{i = 1, \ldots, t} \kappa(\mathfrak q_i)
\quad\text{resp.}\quad
R^{sh} \otimes_R \kappa(\mathfrak p) =
\prod\nolimits_{i = 1, \ldots, s} \kappa(\mathfrak r_i)
$$
where $\mathfrak q_1, \ldots, \mathfrak q_t$,
resp.\ $\mathfrak r_1, \ldots, \mathfrak r_s$
are the prime of $R^h$, resp.\ $R^{sh}$ lying over $\mathfrak p$.
Moreover, the field extensions
$\kappa(\mathfrak q_i)/\kappa(\mathfrak p)$
resp.\ $\kappa(\mathfrak r_i)/\kappa(\mathfrak p)$
are separable algebraic.
\end{lemma}

\begin{proof}
This can be deduced from the more general
Lemma \ref{lemma-filtered-colimit-etale-noetherian-fibres}
using that the henselization and strict henselization are Noetherian
(as we've seen above). But we also give a direct proof as follows.

\medskip\noindent
We will use without further mention the results of
Lemmas \ref{lemma-dumb-properties-henselization} and
\ref{lemma-henselization-noetherian}.
Note that $R^h/\mathfrak pR^h$, resp.\ $R^{sh}/\mathfrak pR^{sh}$
is the henselization, resp.\ strict henselization of $R/\mathfrak p$,
see Algebra, Lemma \ref{algebra-lemma-quotient-henselization}
resp.\ Algebra, Lemma \ref{algebra-lemma-quotient-strict-henselization}.
Hence we may replace $R$ by $R/\mathfrak p$ and assume that $R$
is a Noetherian local domain and that $\mathfrak p = (0)$.
Since $R^h$, resp.\ $R^{sh}$ is Noetherian, it has finitely many
minimal primes $\mathfrak q_1, \ldots, \mathfrak q_t$,
resp.\ $\mathfrak r_1, \ldots, \mathfrak r_s$.
Since $R \to R^h$, resp.\ $R \to R^{sh}$ is flat these are exactly
the primes lying over $\mathfrak p = (0)$ (by going down).
Finally, as $R$ is a domain, we see that $R^h$, resp.\ $R^{sh}$
is reduced, see Lemma \ref{lemma-henselization-reduced}.
Thus we see that $R^h \otimes_R \kappa(\mathfrak p)$
resp.\ $R^{sh} \otimes_R \kappa(\mathfrak p)$
is a reduced Noetherian ring with finitely many primes, all of which
are minimal (and hence maximal). Thus these rings are Artinian and are
products of their localizations at maximal ideals, each necessarily a field
(see Algebra, Proposition \ref{algebra-proposition-dimension-zero-ring} and
Algebra, Lemma \ref{algebra-lemma-minimal-prime-reduced-ring}).

\medskip\noindent
The final statement follows from the fact that $R \to R^h$,
resp.\ $R \to R^{sh}$ is a colimit of \'etale ring maps and hence
the induced residue field extensions are colimits of finite separable
extensions, see
Algebra, Lemma \ref{algebra-lemma-etale-at-prime}.
\end{proof}

\end{document}
